\section{The model and its simultaneous realization}
\label{sam:section}

The expanding incidences and transverse projections now give a model
operator with a separated upper band and small unnormalized trace.
We then realize its compression to \(Y\) by edge voltages and matrix
weights, while imposing the girth needed for the geometric construction.
All norms in this Section, unless a finite compression is explicitly
specified, are operator norms in the full regular representation.  In
particular, the estimates below do not merely test some finite-dimensional
representations of the deck group.

\subsection{The upper space of the model}

Use the notation and parameter choices of the preceding Sections.  The
vertex space before deletion, with its auxiliary indices, is
\[
 \mathcal H^0=
 \ell^2\!\left(\coprod_{i\in\mathbb Z}(A_i\sqcup B_i)\right)
 \otimes\ell^2(F)\otimes\mathbb C^N.
\]
Let \(\Pi\) be the orthogonal projection retaining precisely the vertices
of \(Y\).  Let \(U^0_{ik}:\ell^2(B_k)\to\ell^2(A_i)\) be the
normalized \(Y^0\) incidence from Proposition~\ref{prune:tower-properties}.
It and its adjoint preserve normalized constants, and its norm between
their orthogonal complements is at most \(1-\chi\). Decrease \(\chi\),
if necessary, so that \(0<\chi\le1/4\). Write \(\beta_\ell\) for the
supremum of the deleted fractions in the individual \(A_i,B_i\), so
\(\beta_\ell\le\varepsilon_\ell\to0\) by
Equation~\eqref{prune:total-budget}.
Equation~\eqref{nil:pt} and Lemma~\ref{nil:window} provide the transverse
projections \(p(t)\) and nonnegative layer weights \(h_i(t)\), with
\begin{equation}
 \sum_i h_i(t)^2=1,\qquad
 J(t):=\{i:h_i(t)\ne0\}\text{ consists of at most two adjacent integers}.
 \label{spec:active-layers}
\end{equation}
These assertions concern \(0<t<\delta\). The identity trace of \(p(t)\)
is finite and strictly positive there by
Equation~\eqref{nil:two-factor-trace}. All identity traces include the
ordinary, unnormalized sum over the \(N\) matrix indices.

In the central Fourier decomposition define a self-adjoint
operator field \(O(t)\) by its blocks from \(B_k\) to \(A_i\):
\begin{equation}\label{sam:model}
 O(t)_{A_iB_k}=U^0_{ik}\otimes h_i(t)h_k(t)p(t),
 \qquad |i-k|\le1.
\end{equation}
The blocks in the opposite direction are the adjoints; all remaining
blocks vanish.  The normalized incidence \(U^0_{ik}\) has norm one.
Thus the block Schur estimate bounds \(O(t)\) uniformly, and it defines
an operator on \(\mathcal H^0\).  Here and below the formula is zero off
\(0<t<\delta\).  No continuity at the endpoints is needed for defining
this bounded measurable field.

Let $a_i=|A_i|^{-1/2}\mathbf1_{A_i}$ and
$b_i=|B_i|^{-1/2}\mathbf1_{B_i}$. Denote the transverse fiber Hilbert space,
including its $N$ matrix indices, by $\mathcal H_t$, and put
\[
 \mathcal H^0_t=
 \left(\bigoplus_i\ell^2(A_i)\oplus\bigoplus_i\ell^2(B_i)\right)
       \otimes\mathcal H_t.
\]

\begin{lemma}[The model band and its trace]
\label{spec:model}
For $0<t<\delta$, define an isometry on $p(t)\mathcal H_t$ by
\begin{equation}
 V_*(t)\zeta=2^{-1/2}
 \left(\bigoplus_i h_i(t)a_i\otimes\zeta\right)
 \oplus
 2^{-1/2}\left(\bigoplus_i h_i(t)b_i\otimes\zeta\right),
 \qquad E_*(t)=V_*(t)V_*(t)^*.
 \label{spec:model-isometry}
\end{equation}
Then $O(t)E_*(t)=E_*(t)$, $\|O(t)\|\leq1$, and
\begin{equation}
 O(t)|_{(1-E_*(t))\mathcal H^0_t}\leq(1-\chi)I,
 \qquad
 \operatorname{tr}_{\mathrm{id}}E_*(t)
       =\operatorname{tr}_{\mathrm{id}}p(t).
 \label{spec:model-gap-trace}
\end{equation}
Set $E_*(t)=0$ outside the window; the same upper bound holds there.
Furthermore
\begin{equation}
       \|(1-\Pi)E_*(t)\|\leq\sqrt{\beta_\ell}
 \label{spec:model-cut}
\end{equation}
uniformly in $t$.
\end{lemma}

\begin{proof}
The direct sum of the layerwise constant spaces reduces $O(t)$, since each
$U^0_{ik}$ and its adjoint preserve the corresponding normalized constants.
On these spaces, identify both sides with $\ell^2(\mathbb Z)\otimes
\mathcal H_t$.  Every pair of active indices is an allowed pair by
Equation~\eqref{spec:active-layers}.  The upper-right block is therefore
\[
       |h(t)\rangle\langle h(t)|\otimes p(t),
       \qquad \|h(t)\|_2=1.
\]
The two-sided self-adjoint block has eigenvalue $1$ on the symmetric copy
in Equation~\eqref{spec:model-isometry}, eigenvalue $-1$ on the
antisymmetric copy, and eigenvalue zero on their orthogonal complement.

On the direct sum of the mean-zero spaces, let $\xi=(\xi_k)_k$ lie on the
$B$ side, with transverse coordinates included.  Tensoring a bounded
operator with the projection $p(t)$ does not increase its norm.  Thus the
upper-right block $C(t)$ satisfies
\[
 \|(C(t)\xi)_i\|
 \leq (1-\chi)h_i(t)\sum_k h_k(t)\|\xi_k\|.
\]
Squaring, summing over $i$, and applying the scalar Cauchy--Schwarz
inequality gives $\|C(t)\xi\|\leq(1-\chi)\|\xi\|$.  The self-adjoint
two-sided block consequently has norm at most $1-\chi$ on these spaces.
This proves the norm and spectral assertions.

At a vertex in $A_i$, the diagonal of $E_*(t)$ is
$h_i(t)^2p(t)/(2|A_i|)$, and the corresponding formula on $B_i$ is
$h_i(t)^2p(t)/(2|B_i|)$.  Summing the identity diagonals over both sides and
all layers gives
\[
 \operatorname{tr}_{\mathrm{id}}E_*(t)
 =\left(\frac12\sum_i h_i(t)^2+\frac12\sum_i h_i(t)^2\right)
        \operatorname{tr}_{\mathrm{id}}p(t)
 =\operatorname{tr}_{\mathrm{id}}p(t).
\]
The antisymmetric copy is at $-1$ and makes no contribution to this upper
trace.  Finally, for $\zeta\in p(t)\mathcal H_t$, the squared norm removed
from $V_*(t)\zeta$ is at most
$\frac12\sum_i h_i(t)^2(\beta_\ell+\beta_\ell)\|\zeta\|^2$.
This proves Equation~\eqref{spec:model-cut}.
\end{proof}

The upper space just identified is supported on at most two adjacent
layers, as shown in Figure~\ref{fig:moving-band}. Its vertex factor is a
normalized combination of the constant vectors, so its trace equals the
fiber trace rather than that trace multiplied by the number of vertices.
We must now approximate the compressed model by an actual voltage-cover
operator with finitely many matrix weights, while ensuring the cover's
girth. Section~\ref{spec:section} will use the model gap and deletion
bound to preserve its upper band through this approximation.

\begin{figure}[htbp]
\centering
\begin{tikzpicture}[>=Latex,scale=0.94,every node/.style={font=\small}]
\fill[blue!7,rounded corners] (2.55,-0.5) rectangle (7.45,2.1);
\node[anchor=east] at (-0.2,1.4) {$A$-constants};
\node[anchor=east] at (-0.2,0) {$B$-constants};
\foreach \idx/\x/\lab in {0/1/{i-1},1/3.7/{i},2/6.3/{i+1},3/9/{i+2}} {
 \node[circle,draw,fill=white,inner sep=3pt] (a\idx) at (\x,1.4) {$a_{\lab}$};
 \node[circle,draw,fill=white,inner sep=3pt] (b\idx) at (\x,0) {$b_{\lab}$};
 \node at (\x,-0.85) {layer $\lab$};
}
\draw[blue!70!black,thick] (a1)--(b1);
\draw[blue!70!black,thick] (a1)--(b2);
\draw[blue!70!black,thick] (a2)--(b1);
\draw[blue!70!black,thick] (a2)--(b2);
\node[blue!70!black] at (5,2.45) {$J(t)=\{i,i+1\}$};
\draw[<-,gray] (0.5,-1.7)--(3.7,-1.7);
\draw[->,gray] (6.3,-1.7)--(9.5,-1.7);
\node[align=center] at (5,-1.7) {active layers};
\node at (1.7,-2.2) {$t\downarrow0:\ b(t)\to-\infty$};
\node at (8.3,-2.2) {$t\uparrow\delta:\ b(t)\to+\infty$};
\end{tikzpicture}

\caption{The constant-vector part of the model at a parameter with two
active layers. The blue connections represent the four nonzero blocks
\(h_j(t)h_k(t)p(t)\), not individual graph edges. Their symmetric upper
space is given by Equation~\eqref{spec:model-isometry}. As the parameter
approaches either endpoint, the active layers escape in the indicated
direction and the fiber trace tends to zero at the weighted rate of
Proposition~\ref{nil:weighted}.}
\label{fig:moving-band}
\end{figure}

\subsection{Realizing the model on a voltage cover}

\begin{theorem}\label{sam:realization}
For every fixed \(\eta>0\), and for all sufficiently large \(\ell\),
there are voltages \(z_e\in F\) and matrices \(W_e\in M_N(\mathbb C)\)
on the positively oriented edges of \(Y\) with the following properties.
\begin{enumerate}
\item There is a constant \(C_\eta\), independent of \(\ell\) and the
layer, such that \(\|W_e\|\le C_\eta\).  Augmenting an original tower
label by its matrix weight gives a finite positive alphabet of size
\(o(dD)\).  The normalizing number \(D_{ik}\) is determined by the
original part of the label.
\item The voltage cover has no nonempty cyclically nonbacktracking
closed walk based in layer \(i\) of length at most
\(C_*H_i/\log_2d\).
\item Let \(T_P\), extended by zero on the deleted vertex indices, have
entries from \(B_k\) to \(A_i\)
\begin{equation}\label{sam:entries}
 (T_P)_{ab}=
 \frac1{\sqrt{dD_{ik}}}\sum_{e:a\to b}W_e\rho(z_e),
 \qquad a\in A_i,\quad b\in B_k,
\end{equation}
with the adjoint entries in the opposite direction.  Then
\begin{equation}\label{sam:approx}
 \|T_P-\Pi O\Pi\|\le\eta.
\end{equation}
Both operators have layer width one.  In particular, their central
Fourier fields satisfy the same inequality almost everywhere, including
\(\|T_P(t)\|\le\eta\) outside the open window.
\end{enumerate}
The matrices \(W_e\) are the actual weights of the individual edges of
the voltage cover, so \(T_P\) reduces on each component of that cover.
\end{theorem}

We first reduce regular-norm estimates to finite tests. Finite empirical
lists then restrict the possible matrix weights, while a block-specific
distribution approximates the model mean and retains wide independent
voltage coordinates. After fixing those lists, we draw once on each
edge. The girth and centered-moment estimates apply to these same
independent draws.

\subsection{Finite tests for the full regular norm}

\begin{lemma}\label{sam:boxes}
Let \(x\ge2\).  Suppose a finite set of elements of \(F\), including
their inverses, has every integral coordinate bounded by \(x^a\), where
\(a\) is fixed.  For every fixed \(B>0\) there is a finite coordinate
box \(\mathcal B_x\subset F\) with
\[
 |\mathcal B_x|\le x^{c},\qquad
 \frac{|\mathcal B_x\mathbin\triangle\mathcal B_x g|}
 {|\mathcal B_x|}\le Cx^{-B}
\]
for all the specified elements \(g\).  Constants \(c,C\) depend on
\(a,B,F\), but not on the number of specified elements or on \(x\).

Let \(A\) be a self-adjoint finite matrix of right convolution sums on
\(\ell^2(J)\otimes\ell^2(F)\otimes\mathbb C^n\), where \(J\) is a
finite index set, with these group supports.  Write its coefficients as
\(C_{uv}(g)\in M_n(\mathbb C)\), and suppose
\[
 \max_{u\in J}\sum_{v\in J}\sum_g\|C_{uv}(g)\|\le R.
\]
If \(P_x\) retains the deck indices in \(\mathcal B_x\) and all the
other indices, then
\begin{equation}\label{sam:finite-test}
 \|P_xAP_x\|\le\|A\|
 \le\|P_xAP_x\|+Cx^{-B}R.
\end{equation}
The same conclusion holds for a rectangular convolution matrix after
self-adjoint off-diagonal doubling, using the larger of its absolute
row and column sum bounds.  The constants do not depend on the matrix
amplification \(n\).
\end{lemma}

\begin{proof}
The integral coordinates and their positive integer weights are those
of Lemma~\ref{nil:grading}.  Let a coordinate of weight \(j\) range
between \(-\lfloor T^j\rfloor\) and \(\lfloor T^j\rfloor\), and denote
the resulting box by \(\mathcal B(T)\).  If \(f\in\mathcal B(T)\),
the weight-\(j\) coordinate of \(fg\) minus that of \(f\) is a sum of
finitely many graded monomials.  Every monomial contains a coordinate
of \(g\), and consequently the total weight contributed by coordinates
of \(f\) is at most \(j-1\).  The same observation includes the pure
\(g\) term.  Since the multiplication polynomials are fixed, there are
constants \(P,C\) such that, when \(T\ge x\), all these changes have
absolute value at most \(Cx^P T^{j-1}\).

It follows that a point whose coordinate of every weight \(j\) is at
distance more than \(Cx^P T^{j-1}\) from the two corresponding faces
stays in the box after right multiplication by \(g\).  The fraction
of points in the union of these boundary strips is at most
\(C'x^P/T\); the rounding of the side lengths changes only the
constant.  Right multiplication is a bijection, so the same bound,
with a factor two, controls the relative symmetric difference.  Take
\(T=\lceil x^{P+B+2}\rceil\).  The number of coordinates and their
weights are fixed, so \(|\mathcal B(T)|\le x^c\) after enlarging a
fixed exponent.  This proves the first assertion, uniformly for all
the possible supports.

For the norm assertion, let \(P_s\) retain deck indices in the left
translate \(s\mathcal B_x\).  Right convolution commutes with left
translation, and therefore all \(P_sAP_s\), acting on their ranges,
have the same norm as \(P_xAP_x\).  For a finitely supported vector
\(\xi\),
\[
 \sum_{s\in F}\|P_s\xi\|^2=|\mathcal B_x|\,\|\xi\|^2.
\]
Indeed any fixed deck index belongs to exactly \(|\mathcal B_x|\)
left translates.  Averaging \(\langle AP_s\xi,P_s\xi\rangle\)
over \(s\), and dividing by \(|\mathcal B_x|\), multiplies the
coefficient of \(\rho(g)\) by
\[
 \alpha_x(g)=\frac{|\mathcal B_x\cap\mathcal B_x g^{-1}|}
 {|\mathcal B_x|}.
\]
Call the resulting self-adjoint convolution matrix \(A_x\).  The
averaged quadratic form satisfies
\[
 |\langle A_x\xi,\xi\rangle|
 \le\|P_xAP_x\|\,\|\xi\|^2.
\]
Moreover, \(|1-\alpha_x(g)|\le Cx^{-B}\).  Self-adjointness gives
the same coefficient norm bound for columns as for rows.  The scalar
Schur estimate, applied to the norms of the matrix entries, gives
\(\|A-A_x\|\le Cx^{-B}R\).  Density and the self-adjoint quadratic
form formula now prove \eqref{sam:finite-test}.  No coordinate within
\(\mathbb C^n\) was estimated separately; this explains the asserted
uniformity in amplification.  For a rectangular \(S\), apply the
argument to
\(\left(\begin{smallmatrix}0&S\\S^*&0\end{smallmatrix}\right)\),
whose norm equals \(\|S\|\), before and after the indicated
compressions.
\end{proof}

\begin{lemma}\label{sam:empirical}
Set \(M=\lfloor d^{1/32}\rfloor\), with \(q=d^{A_*}\) and \(A_*\)
fixed.  Consider at most \(C(1+\log q)\) finite distributions of
operators \(W\rho(g)\) on \(\ell^2(F)\otimes\mathbb C^N\).
Suppose \(\|W\|\le B_0\), every possible group coordinate is
bounded by \(q^a\), and \(N\le q^a\), for fixed constants
\(a,B_0,C\).  For each fixed \(\varepsilon>0\), and sufficiently
large \(\ell\), each expectation has an \(M\)-term empirical
average within \(\varepsilon\) in the full regular norm.  The
averages may be chosen for all the distributions simultaneously.
\end{lemma}

\begin{proof}
For each distribution take \(M\) independent samples.  Apply
Lemma~\ref{sam:boxes} at \(x=q\) to the difference of their average
and their expectation, after off-diagonal doubling.  Its absolute
coefficient row and column sums are bounded by \(2B_0\).  Choose the
box to make the error less than \(\varepsilon/2\).  The resulting
self-adjoint test matrix has dimension at most \(q^{c_0}\), with
\(c_0\) fixed.  If \(Z_j\) is its compressed, centered \(j\)-th
sample, then \(\mathbb E Z_j=0\) and \(\|Z_j\|\le2B_0\).

Expand the trace of \((M^{-1}\sum_j Z_j)^{2k}\).  If an index occurs
exactly once in a term, conditioning on all other samples makes that
term zero.  This argument is valid for a noncommutative product: the
unique occurrence has the form \(LZ_jR\) with \(L,R\) fixed under
the conditioning.  At most \(k\) different indices occur in a
remaining term.  The number of sequences using at most \(k\)
indices is at most
\(\sum_{a=1}^k M^a a^{2k}\).  Bounding the absolute trace of each
product by its dimension times the product of its factor norms gives
\[
 \mathbb E\operatorname{Tr}
 \left(M^{-1}\sum_{j=1}^M Z_j\right)^{2k}
 \le q^{c_0}(2B_0)^{2k}M^{-2k}
       \sum_{a=1}^k M^a a^{2k}
 \le q^{c_0}k\left(\frac{4B_0^2k^2}{M}\right)^k.
\]
For \(k=\lceil\log q\rceil\), Markov's inequality therefore bounds
the probability of test norm greater than \(\varepsilon/2\) by
\[
 q^{c_0}k\left(\frac{16B_0^2k^2}{M\varepsilon^2}\right)^k.
\]
Its logarithm is at most
\(-c(\log q)^2+O(\log q\log\log q)\), since
\(\log M=(32A_*)^{-1}\log q+O(1)\).  The bound tends to zero even
after multiplication by \(C(1+\log q)\).  Thus all the tests succeed
with positive probability.  Lemma~\ref{sam:boxes} transfers their
bounds to the full regular norm.  Fix one such list for each
distribution.
\end{proof}

\subsection{Compact coefficient templates and masks}

Fix a small accuracy \(\varepsilon>0\), to be chosen in terms of
\(\eta\) at the end.  In this Subsection all auxiliary errors can be
made smaller than any prescribed fixed multiple of \(\varepsilon\).
The constants they require are fixed before increasing \(\ell\).

By Lemma~\ref{nil:templates}, choose smooth compactly supported functions
\(J_\epsilon^\flat\) and a fixed relative logarithmic mesh size
\(0<u\le1\) so that truncation and coefficient rounding have arbitrarily
small prescribed \(\ell^1\) errors.  These choices depend only on
\(\varepsilon\) and the previously fixed data.
The estimates apply uniformly on the active window because
\(t,\delta-t\in(0,\delta)\) and \(\delta\to0\) as \(\ell\to\infty\);
they give operator estimates for every unitary family used below.
Throughout rounding, the translating operators, and consequently their
multipliers, keep their actual parameter; only the coefficients are rounded.

Choose a sufficiently fine fixed log-grid spacing.  A nonnegative
smooth bump on the real line whose translates at that spacing cover
the line, divided by the positive sum of those translates, gives a
smooth partition of unity.  Pull one such partition back under
\(\log_2(t/\delta)\), and another under
\(\log_2((\delta-t)/\delta)\).  Denote them by \(\psi_\nu(t)\)
and \(\theta_\omega(t)\), and denote their associated positive
magnitude templates by \(a_\nu,b_\omega\).  On the supports in
question,
\[
 |\log(a_\nu/t)|\le u,\qquad
 |\log(b_\omega/(\delta-t))|\le u.
\]
Both partitions sum to one on \(0<t<\delta\).

For a fixed block \((i,k)\), keep precisely the template pairs whose
supports intersect the support of \(h_i h_k\), and call their set
\(\mathcal T_{ik}\).  Its cardinality is bounded by a constant
\(K_\varepsilon\), uniformly in \(i,k,\ell\).  Indeed
\(h_i h_k\ne0\) implies \(|b(t)-i|<1\); each of
\(\log_2(t/\delta)\) and
\(\log_2((\delta-t)/\delta)\) varies by at most \(2\kappa\)
there.  Only boundedly many members of either fixed-spacing grid can
meet these ranges.  The specified positivity of \(h\) ensures that
the nonzero incidence blocks have some such pairs.

For any twisted representation \(U'\) of \(\bar E_{\mathbb Z}\)
define the finite sum
\[
 \mathcal L_\epsilon(\alpha;U')=
 \sum_{\gamma\in\bar E_{\mathbb Z}}
 \alpha^{n_{\mathrm f}}
 J_\epsilon^\flat(\Delta_{\alpha^{1/m}}\gamma)U'(\gamma).
\]
By Proposition~\ref{nil:projection}, Lemma~\ref{nil:templates}, and the
chosen truncation and rounding accuracies, the block field
\begin{equation}\label{sam:template-field}
 \sum_{(\nu,\omega)\in\mathcal T_{ik}}
 f_{ik\nu\omega}(t)\,
 \mathcal L_+(a_\nu;u_t)
 \mathcal L_-(b_\omega;V_t),
 \qquad
 f_{ik\nu\omega}=h_i h_k\psi_\nu\theta_\omega,
\end{equation}
approximates \(h_i h_kp(t)\) to an arbitrarily small fixed error,
uniformly on the window, for sufficiently large \(\ell\).  The two
factors act on their separate transverse indices.  The error estimate
uses the uniformly bounded coefficient \(\ell^1\) norms, positivity
of the scalar partition, and its sum-one property.  Extend each
scalar mask by zero off the window.

The masks have a useful scale uniform in the layer.  Put
\[
 w_i=q^{-1}2^{-\kappa|i|}.
\]
On the support of \(h_i\), the identities
\(t=\delta 2^{\kappa b(t)}/(1+2^{\kappa b(t)})\) and
\(\delta-t=\delta/(1+2^{\kappa b(t)})\) show that the shorter
distance to an endpoint is comparable to \(w_i\), with constants
depending on \(\kappa\).  They also show that the support has
length at most \(Cw_i\).  The chain rule for the two logarithmic
charts gives, for every fixed integer \(j\ge0\),
\begin{equation}\label{sam:mask-derivative}
 \|f_{ik\nu\omega}^{(j)}\|_\infty\le C_jw_i^{-j},
 \qquad
 \|f_{ik\nu\omega}^{(j)}\|_1\le C_jw_i^{1-j}.
\end{equation}
The compact support of the original bumps makes the zero extension
smooth; each individual support is at positive distance from the
endpoints.  In particular, its Fourier coefficients satisfy
\begin{equation}\label{sam:fourier}
 |\widehat f_{ik\nu\omega}(n)|
 \le C_jw_i(1+|n|w_i)^{-j},\qquad n\in\mathbb Z.
\end{equation}
For \(|n|w_i\le1\) this follows from the \(L^1\) bound.  Otherwise
integrate by parts \(j\) times, use
\eqref{sam:mask-derivative}, and combine the two bounds.
Consequently
\[
 \sum_n|\widehat f_{ik\nu\omega}(n)|\le C,
 \qquad
 \sum_{|n|>L/w_i}|\widehat f_{ik\nu\omega}(n)|
 \le C_j(1+L)^{1-j}.
\]
Choose a fixed sufficiently large \(L\), and truncate at
\(K_i=\lceil L/w_i\rceil\).  These are uniform approximations on
the entire circle.  In particular, the approximating trigonometric
polynomials have small values outside the active supports; exact
vanishing there is neither claimed nor used.

\subsection{Finite matrix lists on the second factor}

The scalar masks vary through infinitely many layers, but the final group
must have a finite alphabet of matrix weights. Using a new matrix list at
every layer would not meet that requirement. The range \(i<I\) uses only finitely many second-factor templates, and hence finitely many empirical lists; in the remaining range a congruence
subgroup factorization replaces all the matrix factors by one common list.
The next estimates justify this replacement in the full regular norm.


In blocks with \(i<I\), every second-factor template satisfies
\begin{equation}\label{sam:ordinary-range}
 \log_2 b_\omega^{-1}
 \le \log_2q+\kappa(I+1)+C.
\end{equation}
This follows from the formula for \(\delta-t\) on the support of
\(h_i\), followed by the fixed log rounding.  Thus there are only
\(O(1+\log q)\) distinct such templates over all these blocks, and
their coordinate supports are bounded by a fixed power of \(q\).
Choose a coordinate rectangle containing each compact template
support, with side radii fixed multiples of
\(b_\omega^{-j/m}\) in weight \(j\).  The rectangle has
\(n_\omega\) lattice points.  Uniform sampling of \(v'\) in it,
with the term
\[
 n_\omega b_\omega^{n_{\mathrm f}}
 J_-^\flat(\Delta_{b_\omega^{1/m}}v')
 R(\sigma(v'))\rho(\sigma_Y(v')),
\]
has expectation equal to the ordinary lifted template.  The scalar
multiplier is bounded uniformly, because
\(n_\omega b_\omega^{n_{\mathrm f}}\le C\).
Lemma~\ref{sam:empirical} gives, for each such template, a fixed
\(M\)-term list approximating that expectation in the full regular
norm.  Its entries retain their original scalar multipliers and
matrices.  Fix all these lists now.

For the remaining blocks a single common matrix list suffices.  Let
\(\Lambda=\Lambda_q\) be the subgroup of \(\bar E_{\mathbb Z}\) whose
coordinates are all divisible by \(q\).  Reduction modulo \(q\)
is a surjective homomorphism by the integral polynomial group law;
its kernel is \(\Lambda\), and
\begin{equation}\label{sam:index}
 [\bar E_{\mathbb Z}:\Lambda]=q^{2n_{\mathrm f}}=:Q_q.
\end{equation}
Choose a set \(\mathcal V_q\) of representatives for the cosets
\(\Lambda v\), with each coordinate in
\(\{0,\ldots,q-1\}\).  Every \(\gamma\) is uniquely \(yv\),
with \(y\in\Lambda\), \(v\in\mathcal V_q\).
Equation~\eqref{nil:congruence-R} gives
\(R(\sigma(y))=1\) for \(y\in\Lambda\): its affine shift is
the identity modulo \(q\), and its modulation exponent is an
integer.  The finite average
\begin{equation}\label{sam:coset-average}
 C_q=Q_q^{-1}\sum_{v\in\mathcal V_q}
 R(\sigma(v))\rho(\sigma_Y(v))
\end{equation}
has an \(M\)-term empirical approximation in the full regular norm,
again by Lemma~\ref{sam:empirical}.  Fix one such list, common to
all \(i\ge I\).

For those blocks replace an ordinary second-factor template by the
product
\begin{equation}\label{sam:subgroup-product}
 \left(
 Q_q\sum_{y\in\Lambda}
 b_\omega^{n_{\mathrm f}}
 J_-^\flat(\Delta_{b_\omega^{1/m}}y)\rho(\sigma_Y(y))
 \right)C_q.
\end{equation}
We next prove the needed comparison with the ordinary template,
including its phase and normalization.

The scale estimates from the window and the choices of parameters
give, uniformly for \(i\ge I\) and active templates,
\begin{equation}\label{sam:second-width}
 b_\omega^{-1/m}/q\ge 2^{2E_0H_i}
\end{equation}
when \(\ell\) is sufficiently large.  Here is the explicit
calculation.  Log rounding contributes a bounded constant, and
\[
 \log_2(b_\omega^{-1/m}/q)
 \ge \frac\kappa m i-\left(A_*-\frac{A_*}{m}\right)\log_2d-C.
\]
Since \(i\ge I=H_0\ge\log_2d\) and \(H_i=H_0+i\le2i\),
the hypothesis
\(\kappa/m>10(E_0+A_*+C_\bullet+1)\) leaves more than the
margin needed for \eqref{sam:second-width}.  It follows that every
side length of the dilated subgroup mesh tends to zero uniformly.
Counting its cells shows that the coefficient \(\ell^1\) norm of
the first parenthesis in \eqref{sam:subgroup-product} is bounded.

In the fiber with \(\tau=t-\delta\), the relevant projective
operators satisfy
\[
 V_t(y)V_t(v)
 =\exp(2\pi\mathrm i\tau c(y,v))V_t(yv).
\]
This is the multiplier identity, not a replacement of \(t\) by a
template in the translating operators.  It can also be checked by
multiplying the actual section lifts and then using the center
parameter \(-\delta\) of \(R\).  As
\(R(\sigma(y))=1\), the first parenthesis in
\eqref{sam:subgroup-product} is precisely its prescribed factor in
this product.  The coefficient at \(yv\), before empirical
replacement of \(C_q\), is therefore
\[
 b_\omega^{n_{\mathrm f}}
 J_-^\flat(\Delta_{b_\omega^{1/m}}y)
 \exp(2\pi\mathrm i\tau c(y,v)).
\]
The factor \(Q_q\) in the subgroup sum has canceled the factor
\(Q_q^{-1}\) in the representative average; no matrix trace
normalization occurs here.

Set \(\beta=b_\omega^{1/m}\).  On all the compact supports in this
calculation the elements \(\Delta_\beta v\), with
\(v\in\mathcal V_q\), tend to the identity uniformly, since their
coordinates of weight \(j\) have absolute value at most
\(q\beta^j\le q\beta\).  The polynomial group law implies
\(\Delta_\beta(yv)=(\Delta_\beta y)(\Delta_\beta v)\) and that
this differs from \(\Delta_\beta y\) by a quantity tending to
zero uniformly on those compacts.  Homogeneity of the cocycle gives
\[
 \tau c(y,v)=\frac\tau{b_\omega}
 c(\Delta_\beta y,\Delta_\beta v).
\]
On an active template \(|\tau|/b_\omega\) is bounded, while the
last cocycle tends uniformly to zero.  Smoothness of the compact
coefficient function thus gives a uniformly vanishing difference
between the preceding coefficient and
\(b_\omega^{n_{\mathrm f}}
 J_-^\flat(\Delta_\beta(yv))\).
All possibly nonzero differences lie in one fixed enlarged compact
set in the scaled coordinates.  The normalized number of ordinary
lattice points there is bounded.  Summing the coefficient differences
therefore gives \(o(1)\), uniformly in the active blocks and
parameters.  This proves the comparison with the ordinary template
in operator norm.  Finally, replacing \(C_q\) by its empirical list
costs at most its regular-norm error times the bounded coefficient
\(\ell^1\) norm of the subgroup factor.  The same bound holds almost
everywhere in every central fiber.

All the coefficient sums used above have bounded absolute norms even
at parameters where their templates are inactive.  Thus multiplication
by the masks and replacement of those masks by their Fourier
truncations preserve the error estimates globally on the circle.
More explicitly, on a mask support use its active comparison; off
that support the untruncated mask is zero.  The error made by
substituting the truncated mask is bounded everywhere by its Fourier
tail times the coefficient norm bounds.  Sum over at most
\(K_\varepsilon\) pairs.  This controls the inactive leakage as well
as the active approximation.

\subsection{One draw per edge and the finite alphabet}

We now specify a finite distribution of a single pair \((W,z)\) for
each block \((i,k)\).  This will be sampled independently on its
undirected edges, using the positive orientation and taking the
adjoint weight and inverse voltage for the reverse orientation.
The choices in different parts of the prescription are independent
conditional on the template pair.

First choose \((\nu,\omega)\) uniformly in
\(\mathcal T_{ik}\), then choose an integer \(n\) uniformly in
\(\{-K_i,\ldots,K_i\}\).  Use the total central power \(n\)
with scalar multiplier
\[
 |\mathcal T_{ik}|(2K_i+1)\widehat f_{ik\nu\omega}(n).
\]
By \eqref{sam:fourier} this multiplier is bounded uniformly.
These additional central powers carry no matrix modulation.

Choose \(x\in\bar E_{\mathbb Z}\) uniformly in a full coordinate
rectangle containing the support of
\(J_+^\flat(\Delta_{a_\nu^{1/m}}x)\).  In each coordinate of
weight \(j\), its radius is a fixed positive constant times
\(a_\nu^{-j/m}\), rounded up to an integer.  Write \(n_\nu^+\)
for its number of lattice points.  Use \(\sigma_X(x)\), with
scalar multiplier
\[
 n_\nu^+ a_\nu^{n_{\mathrm f}}
 J_+^\flat(\Delta_{a_\nu^{1/m}}x).
\]
It is uniformly bounded by the rectangle cell-volume estimate.

If \(i<I\), choose one of the fixed \(M\) entries in the ordinary
second-factor list for \(b_\omega\), uniformly, and use its section
lift, its matrix \(R(\sigma(v'))\), and its scalar multiplier.
Uniform averaging supplies the list's normalization \(1/M\).

If \(i\ge I\), choose \(y\in\Lambda\) uniformly in the analogous
full subgroup rectangle, and independently choose one entry \(v\)
uniformly from the fixed representative list.  Use
\(\sigma_Y(y)\sigma_Y(v)\), the matrix \(R(\sigma(v))\), and
scalar multiplier
\[
 n_\omega^\Lambda Q_q b_\omega^{n_{\mathrm f}}
 J_-^\flat(\Delta_{b_\omega^{1/m}}y),
\]
where \(n_\omega^\Lambda\) is the subgroup rectangle's number of
points.  Its normalized cell volume is
\(Q_q b_\omega^{n_{\mathrm f}}\).  By
\eqref{sam:second-width}, each scaled side mesh is at most one,
so this scalar multiplier is also bounded uniformly.

Let \(z\) be the product of the prescribed group elements, in the
indicated order.  Let \(c_e\) be the product of the scalar
multipliers, and let \(R_v\) be the one indicated matrix.  The raw
weight is \(c_eR_v\), with \(|c_e|\le C_\varepsilon\).  Fix a finite
grid in the complex disk of radius \(C_\varepsilon+1\), and round
\(c_e\) deterministically to a grid value \(c_e^\flat\) at distance
at most a prescribed \(\varepsilon'\).  Define
\[
 W=c_e^\flat R_v.
\]
Since \(R_v\) is unitary, the rounding changes the expectation
norm by at most \(\varepsilon'\).  The voltage \(z\) is retained
even when the weight is zero.  In particular, the support of the
voltage distribution is not conditioned on a nonzero weight.

All distributions in this prescription are finite.  Linearity of
expectation shows, term by term, that the raw expectation is the
sum of the truncated Fourier masks times the appropriate first and
second coefficient templates.  Uniform sampling of a rectangle
cancels exactly its point-count multiplier, and uniform sampling of
a template pair and a frequency cancels their two count multipliers.
The empirical lists give their specified averages.  The preceding
Subsections and the last rounding consequently show that we can
arrange
\begin{equation}\label{sam:edge-mean}
 \mathop{\mathrm{ess\,sup}}_t
 \bigl\|\mathbb E[W\rho(z)](t)-h_i(t)h_k(t)p(t)\bigr\|
 \le\varepsilon
\end{equation}
uniformly in \((i,k)\).

Only the empirical-list matrices enter the tags.  Their total number
is at most \(C_\varepsilon(1+\log q)M\), from the ordinary lists
and the single representative list.  The final scalar grid has
fixed finite cardinality depending on \(\varepsilon\) alone.
Thus the total number of weight tags has the same stated bound.
The original tower alphabet has size
\(O(d2^r(1+s)^2)\).  As \(2^r=d^{7/4}\), their product is
\[
 O_\varepsilon\bigl(d^{89/32}(1+\log d)^3\bigr)=o(d^3)=o(dD).
\]
In particular, neither the possibly very large group-coordinate
rectangles nor the number of different first-factor templates enlarge
the label alphabet.  Voltages are separate data of the cover.

There is a fixed constant \(C\), independent of \(i,\ell\), such
that every coordinate of every possible \(z\) in a block with first
index \(i\), and of its inverse, is bounded by
\begin{equation}\label{sam:support}
 2^{CH_i}.
\end{equation}
Indeed the logarithms of the template radii and of \(K_i\) are
\(O(\log q+\kappa|i|+1)=O(H_i)\).  The representative coordinates
are bounded by \(q\); the ordinary empirical terms lie in their
specified rectangles.  Fixed polynomial multiplication and inversion
preserve a bound of the form \eqref{sam:support}.

The first factor supplies wide independent coordinates in a range
overlapping that of \eqref{sam:second-width}.  Namely, for every
active first-factor template with \(i\le2I\),
\begin{equation}\label{sam:first-width}
 a_\nu^{-1/m}\ge 2^{2E_0H_i}
\end{equation}
for sufficiently large \(\ell\).  For \(0\le i\le2I\), the
lower bound
\(\log_2a_\nu^{-1}\ge\log_2q-C\), together with
\(H_i\le3H_0\le3h_*\log_2d\) and
\(A_*/m>10E_0h_*\), proves it.  For \(i<0\), use instead
\[
 \log_2a_\nu^{-1}
 \ge\log_2q+\kappa|i|-C
\]
and \(H_i=H_0+|i|\), with the same inequality for \(A_*\) and
the stronger chosen inequality for \(\kappa\).
Every upper unitriangular entry in this first rectangle is sampled
independently on an integer interval containing at least
\(2^{E_0H_i}\) points.  For \(i\ge I\), every corresponding
entry of the second-factor subgroup rectangle has at least that
many possibilities in its step-\(q\) interval.  Fixed rectangle
radii and integer rounding are harmless because
\eqref{sam:first-width} and \eqref{sam:second-width} have an
exponential margin.

\subsection{Girth with the same independent edge draws}

We use the Schwartz--Zippel polynomial zero bound, whose short proof we include.
The sharp total-degree bound is due to Schwartz
\cite[Lemma~1, p.~702, and Corollary~1]{Schwartz1980}; Zippel's sparse-interpolation
work \cite[\S 3.1, Theorem~1]{Zippel1979} gives related coordinate-degree
zero estimates, and DeMillo and Lipton \cite{DeMilloLipton1978} earlier
used random evaluations for polynomial identity testing.
The minimum-side bound below is an immediate consequence of Schwartz's
estimate for unequal coordinate sets. If a nonzero
polynomial of total degree at most \(r\), over a field, is evaluated
on independent uniform choices from finite sets of sizes at least
\(L\), its probability of vanishing is at most \(r/L\).  This
follows by induction on the number of variables.  Regard the
polynomial as a polynomial of degree \(s\) in the last variable;
its leading coefficient is a nonzero polynomial of degree at most
\(r-s\).  The probability that this coefficient vanishes is at most
\((r-s)/L\); otherwise the one-variable root bound gives at most
\(s/L\).  The constant case starts the induction.  Arithmetic
progressions of distinct integers are allowed as the finite sets.

Fix a nonempty cyclically nonbacktracking base walk \(w\), based in
layer \(i\), of length
\(a\le C_*H_i/\log_2d\).  The layer indices along it satisfy
\(|H_j-H_i|\le a+1\), so their ratios to \(H_i\) tend uniformly
to one.  If some positive-edge block encountered has first index
\(j_0<I\), then all such first indices are at most \(2I\) when
\(\ell\) is sufficiently large.  To verify this also for very
negative \(j_0\), the length estimate and
\(H_i\le H_{j_0}+a+1\) give
\(a+1\le\epsilon_\ell H_{j_0}\), with
\(\epsilon_\ell\to0\).  The maximum first index is at most
\(j_0+\epsilon_\ell(I+|j_0|)\).  If \(j_0\ge0\), this is
at most \((1+2\epsilon_\ell)I\); if \(j_0<0\), it is at most
\(\epsilon_\ell I\).  In this case use the first-factor quotient
of \(F\) onto the upper unitriangular group.  If there is no such
block, every first index is at least \(I\), and use the second
upper unitriangular quotient.  These are the two overlapping wide
ranges just proved.

Assign an independent formal upper unitriangular matrix to each
undirected edge in the support of \(w\), using its inverse when
the edge is traversed backwards.  Lemma~\ref{nil:word} says that
the resulting word is not identically the identity.  Its hypotheses
apply because the support has the tower girth and the present walk
has length at most the fixed \(C_*\) scale.  In particular, the
integer \(m\) used for this assertion was fixed at that scale.
Some entry of this word minus the identity is therefore a nonzero
polynomial in the upper entries of those edge matrices.  Its total
degree is at most \(m-1\).  Indeed every term in the \((u,v)\)
entry of a product of upper unitriangular matrices, including inverse
matrices expanded by the finite geometric series, has total
superdiagonal weight \(v-u\).  Each variable has positive integer
superdiagonal weight, so its ordinary total degree is at most
\(v-u\le m-1\).  This bound is independent of the number of
factors in the word.

Condition on the template choices, frequencies, empirical-list
choices, the other transverse factor, and the other independently
sampled coordinates.  Do not condition on the final scalar weights.
In the first-factor case the remaining upper entries have their
independent uniform interval distributions.  In the second-factor
case the edge matrix is a random subgroup upper matrix multiplied
on the right by its fixed representative upper matrix.  Right
multiplication by that fixed matrix is an invertible polynomial
coordinate change over the reals.  Replacing the subgroup entries
by their multiples of \(q\) is likewise an invertible linear
coordinate change over the reals.  Thus the nonzero word polynomial
remains nonzero in the independently sampled entries.  This also
explains why weights that happen to vanish do not remove a voltage
variable from this argument.

Each variable set has at least \(2^{E_0H_j}\) points at its local
layer \(j\).  The polynomial estimate therefore bounds the
conditional, and hence unconditional, probability of trivial
holonomy by
\[
 (m-1)2^{-E_0\min_j H_j}.
\]
There are at most
\[
 (d+1)2^{H_i}
 \sum_{1\le a\le C_*H_i/\log_2d}(6D)^a
 \le 2^{(3C_*+3)H_i}
\]
walks to test from all vertices of layer \(i\), for sufficiently
large \(\ell\).  Here the rough bound counts all walks, and the
maximum degree is at most \(6D\).  Since
\(\min_jH_j\ge(1-o(1))H_i\) and
\(E_0>10(C_*+C_\bullet+1)\), the union bound is at most
\(C2^{-2H_i}\) after increasing \(\ell\).  Summing over all
\(i\in\mathbb Z\) gives a bound tending to zero.

A closed lift projects under a graph covering to a closed walk,
and nonbacktracking, including the cyclic condition, is preserved.
Trivial holonomy depends only on the base walk, not on which lift
of its initial vertex is used.  Consequently the tests just counted
exclude every forbidden closed walk in the voltage cover.

\subsection{Centered moments, with cycle edges included}

To control deviations from the block means, we expand even moments
of the centered edge operators. Independence and centering make a
term vanish whenever its base walk traverses an undirected edge
exactly once. The next lemma counts the remaining walks when their
supports have bounded cycle rank; its constants do not affect the
earlier nilpotence parameter.

\begin{lemma}\label{sam:walk-count}
Let a finite graph have maximum degree \(\Delta\ge1\).  Fix a
vertex \(v\) and an integer \(l\ge1\).  Suppose the support of
every closed walk under consideration has first Betti number at
most \(b\).  The number of length-\(2l\) closed walks from \(v\)
in which no undirected edge occurs exactly once is at most
\[
 C_b^l\Delta^l.
\]
Multiple edges, if present, are treated as distinct edges.
\end{lemma}

\begin{proof}
The support of such a walk is connected and has at most \(l\)
edges.  Prune leaves other than \(v\) successively, and retain the
resulting rooted core \(K\).  If the support is a tree, the core
is the single vertex \(v\).  Otherwise the core contains \(v\),
has no vertex of degree one other than possibly \(v\), and has
the same first Betti number as the support.  Write \(a\) for its
edge count.  The identity
\[
 \sum_{x\in K}(\deg_K(x)-2)=2b_1(K)-2
\]
gives a maximum-degree bound of \(2b+2\). If \(K\) has no edges,
its maximum degree is zero. Otherwise all terms in the sum are
nonnegative except possibly the root term, which is at least \(-1\).
The parts removed by pruning are trees, each
attached at exactly one core vertex.  Otherwise a remaining cycle,
or a path joining two core vertices, could not have been removed
by the leaf procedure.

There are at most \(4^a\Delta^a\) possibilities for a rooted
connected subgraph \(K\) with \(a\) edges.  We give an encoding
that includes its cycle edges.  Starting at \(v\), when an as yet
unmarked incident edge of \(K\) is available, mark it, traverse it,
and push that edge onto a stack.  Explore recursively at its other
endpoint, even if that vertex has been visited before.  When this
recursive call has exhausted its available unmarked edges, return
along the pushed edge and pop it.  Marking \emph{edges}, rather
than vertices, ensures that each edge, including every edge closing
a cycle, is pushed exactly once and popped exactly once.  The
process explores the connected subgraph, and it has a balanced
push/pop word of length \(2a\), at most \(4^a\) possibilities.
At a push, record the chosen incident edge, with at most \(\Delta\)
possibilities.  At a pop the edge is determined by the stack.
Given the starting vertex, this code reconstructs the complete
traversal and hence its edge set.  One may make the encoding
deterministic using fixed orders on incident edges.  This proves
the bound; it does not assume that the explored subgraph is a
tree.

Fix a possible core \(K\).  Encode a walk by declaring each step
to be a core step, a tree descent away from the core, or a tree
return toward the core.  There are at most \(3^{2l}\) such words.
For each core step there are at most \(2b+2\) choices.  A descent
has at most \(\Delta\) choices, while a return is forced by the
stack of the current tree excursion.  Let \(c\) be the number of
core steps and \(t\) the number of descents.  Returns and descents
are equally numerous, since each excursion begins and ends at
its attachment vertex.  Thus \(2t+c=2l\).  Every core edge is
used at least twice, by the no-singleton assumption, whence
\(c\ge2a\) and \(t\le l-a\).

The number for this fixed core is bounded by
\(3^{2l}(2b+2)^{2l}\Delta^{l-a}\).  This is an upper bound even
if some of the formal encodings fail to be valid tree excursions.
Multiply by the core count and sum over \(0\le a\le l\).  Since
\(l+1\le2^l\) and \(4^a\le4^l\), the resulting bound is
\(\bigl(72(2b+2)^2\bigr)^l\Delta^l\), as required.
\end{proof}

Fix a block \((i,k)\) of the actual surviving graph.  Before degree
normalization let \(X_{ik}\) be its centered rectangular convolution
matrix, with entries
\[
 (X_{ik})_{ab}=
 \sum_{e:a\to b}
 \bigl(W_e\rho(z_e)-\mathbb E[W_e\rho(z_e)]\bigr).
\]
Double it off the diagonal to a self-adjoint matrix.  By
\eqref{sam:support}, Lemma~\ref{sam:boxes} applies with
\(x=2^{H_i}\), simultaneously for every possible assignment of
these finitely many edges.  The row and column coefficient bounds
are \(O_\varepsilon(D)\).  Choose the boundary exponent so that
the absolute error in \eqref{sam:finite-test} tends uniformly to
zero.  Its finite test dimension, including all vertex, deck, and
matrix indices, is bounded by
\begin{equation}\label{sam:dimension}
 2^{C_1H_i}
\end{equation}
for a fixed \(C_1\).  To check this, the block has
\(O((d+1)2^{H_i+1})\) vertices, the deck box has
\(2^{O(H_i)}\) points, and
\(N=q^{n_{\mathrm f}}=2^{A_*n_{\mathrm f}\log_2d}
 \le2^{O(H_i)}\).

At this point fix a constant
\begin{equation}\label{sam:C2-order}
 C_2>2(C_1+4),\qquad
 l=\left\lceil C_2H_i/\log_2d\right\rceil.
\end{equation}
The tower girth implies that the support of every length-\(2l\)
walk with no singleton edge has cycle rank bounded by a constant
depending on \(C_2\) alone, using Lemma~\ref{nil:rank}.  Indeed it
has at most \(l\) edges and lies in the two fixed layers \(i,k\),
whose girth lower bounds are comparable to
\(H_i/(100\log_2d)\).  The ratio of the edge count to that lower
bound is bounded in terms of \(C_2\).  This use of
Lemma~\ref{nil:rank} requires no identity test for a longer word.
In particular \(m\), which was fixed for the \(C_*\) girth
tests, is not increased when \(C_2\) is chosen.

Let \(\widehat X_{ik}\) denote the finite self-adjoint compression.
Every centered compressed edge entry has norm at most
\(2C_\varepsilon\).  Expand
\(\mathbb E\operatorname{Tr}(\widehat X_{ik}^{2l})\) along the
base-vertex indices.  A summand whose walk has a singleton
undirected edge has expectation zero: condition on the other
edges, and the sole remaining factor is a centered entry or its
adjoint.  This remains true after deck compression.  For the other
walks, the absolute trace of the product on a fixed vertex's deck
and matrix indices is at most their dimension times
\((2C_\varepsilon)^{2l}\).  Lemma~\ref{sam:walk-count}, with
\(\Delta\le6D\), now gives
\begin{equation}\label{sam:block-moment}
 \mathbb E\operatorname{Tr}(\widehat X_{ik}^{2l})
 \le 2^{C_1H_i}(C_4D)^l,
\end{equation}
where \(C_4\) is fixed after \(C_2\) and the accuracies have
been chosen.  It is independent of \(i,\ell\).  The exponent in
the dimension bound has not changed in this argument.

Use the finite norm test and Markov's inequality at threshold
\(\varepsilon\sqrt{dD_{ik}}\).  For large \(\ell\), the absolute
box error is less than half this threshold.  Since
\(D/2\le D_{ik}\le2D\), \eqref{sam:block-moment} implies
\begin{equation}\label{sam:block-probability}
 \mathbb P\{\|X_{ik}\|>
                   \varepsilon\sqrt{dD_{ik}}\}
 \le 2^{C_1H_i}(C_5/d)^l.
\end{equation}
The constant \(C_5\) can be very large and can depend on \(C_2\).
Nevertheless it is fixed before increasing \(\ell\).  Increase
\(\ell\) so that \(C_5\le d^{1/2}\).  The right side is then
at most
\(2^{(C_1-C_2/2)H_i}\le2^{-4H_i}\), by
\eqref{sam:C2-order}.  Summing over the at most three blocks with
each first index gives a total failure bound tending to zero.

\subsection{A simultaneous choice and the global norm}

For sufficiently large \(\ell\), the sum of the girth failure
bounds and the block-norm failure bounds is less than one.  In the
countable independent product of the finite edge distributions,
the union bound therefore gives positive probability that all
requirements hold.  There is also an elementary compactness
formulation which avoids any need for an infinite sampling
construction.  Enumerate the countably many edge choices and the
countably many requirements.  Every requirement depends on only
finitely many edges: this is immediate for a fixed walk, and a
fixed block contains finitely many edges even though its convolution
norm acts on an infinite deck space.  Any finite collection of
requirements is satisfiable by the finite-product union bound,
using the same total budget below one.  Choose assignments
satisfying the first \(n\) requirements and extend them arbitrarily
to all other edges.  Successive subsequences make the choice at
each edge constant, since every edge has finitely many possible
outcomes.  The diagonal limit satisfies every requirement: once
its finite set of edges has stabilized, the corresponding operators
and voltages are exactly the same.  Fix such an assignment.

It remains to combine the block errors.  Because the distribution
is common to all edges in a block, the expected normalized block
equals the actual surviving normalized incidence tensored with its
single-edge expectation.  The surviving incidence is the compression
of \(U^0_{ik}\) and has norm at most one.  Equation
\eqref{sam:edge-mean} therefore bounds its difference from the
corresponding block of \(\Pi O\Pi\) by \(\varepsilon\).
The simultaneous success of \eqref{sam:block-probability} bounds
the centered normalized block by \(\varepsilon\).  A matrix with
layer width one and block norms bounded by \(\varepsilon\) has
norm at most \(3\varepsilon\): apply Cauchy--Schwarz to each row
and sum, or apply the scalar row-and-column Schur bound.  Its
self-adjoint off-diagonal doubling has the same norm.  Consequently
\[
 \|T_P-\Pi O\Pi\|\le6\varepsilon.
\]
Taking \(\varepsilon<\eta/6\) proves
\eqref{sam:approx}.  All expectation errors used above were made
smaller than this chosen \(\varepsilon\) by first fixing compact
cutoffs, log-grid spacing, empirical accuracies, Fourier cutoffs,
and scalar rounding, and then increasing \(\ell\).  The order of
the constants is thus: the geometric and nilpotence parameters;
the fixed desired accuracy and its coefficient data; the finite
test exponent \(C_1\); the moment constant \(C_2\); and finally
\(\ell\).  All estimates are uniform in the layers.  This proves
Theorem~\ref{sam:realization}.
